
Low-rank adaptation (LoRA) is widely applied to transformer models by targeting \texttt{nn.Linear} projection layers --- a recipe that practitioners assume transfers to any architecture with equivalent layer types. We test this assumption on Mamba-130m, a pure state space model (SSM) whose projection layers are \texttt{nn.Linear} and whose community documentation reports 90--92\% SST-2 accuracy with projection-only LoRA. Applying LoRA to \texttt{in\textbackslash \_proj}, \texttt{out\textbackslash \_proj}, and \texttt{x\textbackslash \_proj} (rank 8, 1.14\% trainable parameters), we find that SST-2 accuracy remains at 0.5092 across all three training epochs --- statistically indistinguishable from the zero-shot baseline of 0.4908 --- while MNLI accuracy degrades from 0.3463 to 0.3234 across three epochs, QNLI oscillates near its zero-shot baseline (0.5046) with no meaningful improvement, and QQP remains at 0.0000. All results fall well below the conservative gate criterion of SST-2 > 0.70. The failure is not a configuration error: LoRA installs correctly and adapter weights receive gradient updates, but those updates carry no task-discriminative information. Notably, experiments ran under the sequential Python fallback implementation of Mamba (the custom CUDA scan kernel was not installed in the experimental environment), confirming that the gradient barrier effect does not depend exclusively on the CUDA kernel implementation. These findings establish that API compatibility with \texttt{nn.Linear} is not a sufficient condition for LoRA to produce learning on Mamba, and that Mamba parameter-efficient fine-tuning requires gradient-path-aware methods that operate inside or upstream of the SSM scan.

